% -----------------------------------------------------------
\section{Honor, Honorable}
% -----------------------------------------------------------
	\label{HONOR}
	
% ----------------------
\subsection{See also}
% ----------------------

	\hyperref[PRIESTHOOD]{Priesthood},

% ----------------------
\subsection{1828 American Dictionary of the English Language} % *1828
% ----------------------
	\label{HONOR-1828}
	
	% -----
	\subsubsection{Noun}
	% -----

		\begin{compactitem}
			\item[1] The esteem due or paid to worth; high estimation.
			\item[2] A testimony of esteem; any expression of respect or of high estimation by words or actions; as the \textit{honors} of war; military \textit{honors}; funeral \textit{honors}; civil \textit{honors}.
			\item[3] Dignity; exalted rank or place; distinction.
			\item[4] Reverence; veneration; or any act by which reverence and submission are expressed, as worship paid to the Supreme Being.
			\item[5] Reputation; good name; as, his \textit{honor} is unsullied.
			\item[6] True nobleness of mind; magnanimity; dignified respect for character, springing from probity, principle or moral rectitude; a distinguishing trait in the character of good men.
			\item[7] An assumed appearance of nobleness; scorn of meanness, springing from the fear of reproach, without regard to principle; as, shall I violate my trust? Forbid it, \textit{honor}.
			\item[8] Any particular virtue much valued; as bravery in men, and chastity in females.
			\item[9] Dignity of mien; noble appearance.
			\item[10] That which \textit{honors}; he or that which confers dignity; as, the chancellor is an \textit{honor} to his profession.
			\item[11] Privileges of rank or birth; in the plural.
			\item[12] Civilities paid.
			\item[13] That which adorns; ornament; decoration.
			\item[14] A noble kind of seignory or lordship, held of the king in capite.
		\end{compactitem}
		
	% -----
	\subsubsection{Verb}
	% -----

		\begin{compactitem}
			\item[1] To revere; to respect; to treat with deference and submission, and perform relative duties to.
			\item[2] To reverence; to manifest the highest veneration for, in words and actions; to entertain the most exalted thoughts of; to worship; to adore.
			\item[3] To dignify; to raise to distinction or notice; to elevate in rank or station; to exalt. Men are sometimes \textit{honor}ed with titles and offices, which they do not merit.
			\item[4] To glorify; to render illustrious.
			\item[5] To treat with due civility and respect in the ordinary intercourse of life. The troops \textit{honor}ed the governor with a salute.
			\item[6] In commerce, to accept and pay when due; as, to \textit{honor} a bill of exchange.
		\end{compactitem}

% % ----------------------
% \subsection{Oxford English Dictionary} % *OED
% % ----------------------
	% \label{HONOR-OED}

	% \begin{compactitem}
		% \item[]
	% \end{compactitem}
	
% ----------------------
\subsection{Buck's Theological Dictionary (1833)} % *BTD
% ----------------------
	\label{HONOR-BTD}

	\begin{compactdesc}
		\item[462--427] HONOUR, a testimony of esteem or submission, expressed by words and an exterior behaviour, by which we make know the veneration and respect we entertain for any one, on account of his dignity or merit. The word is also used in general for the esteem due to virtue, glory, reputation, and probity; as also, for an exactness in performing whatever we have promised; and in this last sense we use the term, \textit{a man of honour}...In order, then to discern where true honour lies, we must not look to any adventitious circumstance, not to any single sparkling quality, but to the whole of what forms a man; in a word, we must look to the soul. It will discover itself by a mind superior to fear, to selfish interest, and corruption; by an ardent love to the Supreme Being, and by a principle of uniform rectitude. It will make us neither afraid nor ashamed to discharge our duty, as it relates both to God and man. It will influence us to be magnanimous without being proud; humble without being mean; just without being harsh; simple in our manners, but manly in our feelings.
			\begin{compactitem}
				\item The second part of this quotation, after the ellipse, comes from some other source, but it isn't cited clearly.
			\end{compactitem}
	\end{compactdesc}
	
% % ----------------------
% \subsection{Restoration Lexicon} % *RL *RESTORATION LEXICON
% % ----------------------
	% \label{HONOR-OED}

	% \begin{compactitem}
		% \item[]
	% \end{compactitem}

% ----------------------
\subsection{Scriptural Usage} % *SCRIPTURES
% ----------------------
	\label{HONOR-SCRIPTURES}

	% % -----
	% \subsubsection{Old Testament} % *OT
	% % -----
		% \label{HONOR-OT}
	
		% % --
		% \paragraph{KJV}
		% % --
			% \label{HONOR-OT-KJV}
	
			% \begin{tabular}{ll}
				% Total occurrences in the OT & 
			% \end{tabular}
			
			% \begin{compactdesc}
				% \item[]
			% \end{compactdesc}
			
		% % --
		% \paragraph{Hebrew Bible}
		% % --
			% \label{HONOR-HEBREW}
		
			% %
			% \subparagraph{\texorpdfstring{\he{sample word}}{}} % paste actual hebrew text here
			% %
				% \label{PAGA} % whatever is in step bible without periods
			
				% \begin{compactdesc}
					% \item[HALOT , PDF ] \textbf{} % Use the 1994 PDF
						% \begin{compactitem}
							% \item[1]
						% \end{compactitem}
					% \item[BDB , PDF ] \textbf{}
						% \begin{compactitem}
							% \item[1]
						% \end{compactitem}
					% \item[DCH PDF ] \textbf{} % don't include the actual page number, they repeat from volume to volume
						% \begin{compactitem}
							% \item[1]
						% \end{compactitem}
				% \end{compactdesc}
				
				% \begin{compactdesc}
					% \item[Some OT uses] \textbf{}
						% \begin{compactdesc}
							% \item[]
						% \end{compactdesc}
				% \end{compactdesc}
			
		% % --
		% \paragraph{Septuagint}
		% % --
			% \label{HONOR-LXX}
		
			% %
			% \subparagraph{\texorpdfstring{\gr{sample word}}{}}
			% %
		
	% -----
	\subsubsection{New Testament} % *NT
	% -----
		\label{HONOR-NT}
	
		% --
		\paragraph{KJV}
		% --
			\label{HONOR-NT-KJV}		
	
			% \begin{tabular}{ll}
				% Total occurrences in the NT & 
			% \end{tabular}
			
			\begin{compactdesc}
				\item[{\hyperref[JOHN-4-44]{John 4:44}}]
				\item[{\hyperref[JOHN-5-23]{John 5:23}}]
				\item[{\hyperref[JOHN-8-49]{John 8:49}}]
				\item[{\hyperref[HEBREWS-5-4]{Hebrews 5:4}}]
			\end{compactdesc}
			
		% --
		\paragraph{Greek New Testament} % *GREEK
		% --
			\label{HONOR-GREEK}
		
			%
			\subparagraph{\texorpdfstring{\gr{τιμή}}{}}
			%
				\label{TIME-2}
			
				\begin{compactdesc}
				
					\item[BDAG 893b] \textbf{}
						\begin{compactitem}
							\item[1] \textbf{the amount at which something is valued, \textit{price, value}}...especially \textit{selling price}...Also concrete \textit{the price received} in selling something
							\item[2] \textbf{manifestation of esteem, \textit{honor, reverence}}
							\item[2.a] active, \textit{the showing of honor, reverence}, or \textit{respect} as an action...usually as a commendation for performance...In \hyperref[1THESSALONIANS-4-4]{1 Thessalonians 4:4} \gr{τιμή} may well be understood in this sense, if \gr{σκεῦος} refers to a female member of the household
							\item[2.b] passive \textit{the respect} that one enjoys, \textit{honor} as a possession
								\begin{compactitem}
									\item There are many examples of this use, along with its combination with other words, including \gr{δόξα}
								\end{compactitem}
							\item[2.c] as a state of being, \textit{respectability}
							\item[2.d] \textit{place of honor}, (\textit{honorable}) \textit{office}...\gr{οὐχ ἑαυτῷ τις λαμβάνει τὴν τιμήν} \textit{no one takes the office of his own accord} \hyperref[HEBREWS-5-4]{Hebrews 5:4}
							\item[3] \textbf{honor conferred through compensation, \textit{honorarium, compensation}}
							\item[4] \textbf{a right that is specially conferred, \textit{privilege}} \hyperref[1PETER-2-7]{1 Peter 2:7}
						\end{compactitem}
						
					\item[CGL 1380a, PDF 1420] \textbf{}
						\begin{compactitem}
							\item[1] \textbf{honour, esteem} (paid by men to gods or superiors, or by gods to men as a reward for services)...plural, honours or marks of honour
							\item[2] (often plural) \textbf{honour, prerogative, privilege} (attaching to a specific position or office, especially that of a god or king)...plural, honoured powers (of underworld gods); honours or privileges (metonymy for those who hold them, i.e. rulers and gods)
							\item[3] \textbf{office} (civic or other)
							\item[4] \textbf{price, value, worth} (of things)...\textbf{return, profit} (from an activity)
							\item[5] \textbf{assessment of value} (of material or abstract things); (referring to wealth) \textbf{criterion} or \textbf{standard of valuation} (with genitive for things)
							\item[6] \textbf{recompense, compensation} (especially for damage to one's personal honour)
						\end{compactitem}
						
					\item[TDNT] \phantom{.}
						\begin{compactdesc}
							\item[8:170--171] \gr{τιµή} has in the first instance a strong material orientation. Odysseus' honour is inseparably bound up with the restoration of control of his possessions, Hom.Od., 1, 117. Achilles' honour is functionally dependent on the number of gifts brought to him to persuade him to take part in the battle, II. 9, 605. Here bodily soundness, the undisputed exercise of social influence and uninfringed enjoyment of one's property are the basis of esteem. Later the noun is used in a more strongly ethical context. A certain type of moral conduct is prerequisite for the esteem a man enjoys. Gradually \gr{τιµή} detaches itself from real possessions and becomes an abstract concept of honour. That the original elements in the meaning of the word were never wholly lost can be seen in the fact that in the \textit{koine} \gr{τιµή} can mean both `honour' and `price.' If in the early Greek period honour as esteem by society on account of concrete circumstances was one of the highest values among the nobility of the 8th cent. B.C. (Hom.), in the city states, esp. Sparta and Athens, the honour of the individual was also that of the \textit{polis}. When under the influence of Sophism the individual came to be increasingly detached from the \textit{polis} the concept of honour became much more individualistic, esp. in Isocrates. But Plato was the first to establish the personal ethical element in honour, or `inward honour,' though without absolutely rejecting `outward honour' (the distinctions accorded a man by the world around). In relation to this wise moderation is to be commended. Plato, then, finally anchored honour in the moral person. The most significant attempt to provide a scientifically grounded ethics of honour was that made by Aristot. The discussion in Eth. Nic., IV, 7, p. 1123b, in which he speaks of \gr{µεγαλοψυχία}, is basic here. The high-minded man must be virtuous, for there is no honour without virtue. He thus possesses honour on the basis of inner worth. By reason of his \gr{ἀρετή} honour is then shown him from without, by his fellow-citizens. If at bottom the high-minded man can only give himself the honour worthy of his virtue, he is in the last analysis above `outward honour.' But there is no honour worthy of perfect virtue. In the Aristot. concept of honour, there is thus a strong individualistic tendency, though the solidarity of the \textit{polis} is not destroyed, for man is by nature a creature destined for political society. Finally Stoicism brought the individualistic concept of honour to its full development. In it `inward honour,' the sense of one's own worth, is decisive. Stoic philosophy was not against every kind of outward honour, but the wise man is relaxed in relation to it; he does not chase it and can do without it. This attitude corresponds to the inner freedom which rules his thought. From the various standpoints the teaching of honour was of great importance among the Greeks and Romans.
							\item[8:175] The combining of \gr{τιµή} and \gr{δόξα}, which is a familiar one in Hellenistic thought and usage ($\rightarrow$ 172, 35 f.; 173, 21 ff., 40), occurs in christological statements in 2 Pt. and Hb. 2 Pt. 1:17 speaks of the transfiguration of Jesus. While the two terms do not occur in Mk. 9:2--8 and Mt. 17:1--8---Lk. alone has \gr{ἐν δόξῃ} in 9:31---a decisive statement in 2 Pt. is that at the transfiguration on the mount Jesus received honour and glory from the Father. \gr{τιµή} and \gr{δόξα} are associated in Hb. 2:7, 9 and 3:3. Hb. 2:6--8a is a quotation from Ps. 8:5--7 LXX, where man's position in creation is denoted by the words \gr{δόξα} and \gr{τιµή}. But that which according to the wording of the Ps. must be applied to man generally or to the individual (\gr{υἱὸς ἀνθρώπου}) is related by the author of Hb. to Christ, 2:9 ($\rightarrow$ III, 1089, 26 ff.). Christ's passion is the presupposition of His crowning with glory and honour. This crowning is obviously His institution into the high-priestly office. \hl{The terms} \gr{δόξα} \hl{and} \gr{τιµή} \hl{thus refer not so much to the glorifying and exalting of Christ as to His high-priestly dignity. The reference might well be to the position of honour which He assumes as the heavenly High-priest.} In Hb. 3:3 Christ and Moses are compared in respect of their \gr{δόξα} and \gr{τιµή}. As the Son of God Christ is far superior to the representative of the old covenant. Certainly the position of Moses was most honourable. To him was entrusted the guiding of the ``house'' of Israel, and of his faithfulness as a servant of God there is no doubt. But since he himself was only part of the house, he is far behind its founder in glory (dignity), $\rightarrow$ V, 125, 16 ff. 30 The superiority of Christ over Moses may be deduced from the very wording of Nu. 12:7. When \gr{δόξα} and \gr{τιµή} are associated, \gr{δόξα} is the higher term. \gr{τιµή} in the sense of a position of honour constitutes only one part of \gr{δόξα}. \hl{In Hb. 5:4} \gr{τιµή} \hl{has very plainly the sense of ``dignity of office.'' Aaron did not take his high-priestly ``dignity'' to himself but received it by divine calling, Ex. 28:1 f. In the same way Christ was instituted in the high-priestly office by God ($\rightarrow$ III, 275, 19 ff.; 277, 17 ff.). But this is an eternal priesthood after the order of Melchisedec.}
						\end{compactdesc}
						
					\item[TDNTA] \phantom{.}
						\begin{compactdesc}
							\item[1072] \gr{τιμή} has at first a strong material orientation to possessions, strength, or social influence. Later, moral conduct plays a bigger part. The fact that \gr{τιμή} can also mean ``price'' upholds the material connection, but \gr{τιμή} as honor increasingly becomes inner worth as distinct from outward esteem. For Aristotle there is no honor without virtue; only on the basis of virtue should outward honor be shown. No honor is enough for perfect virtue, and the person of inner worth is finally above outward honor. In Stoicism inner honor is what counts. The sage, enjoying inward freedom, can live without external honor and can thus be relaxed in relation to it.
						\end{compactdesc}
						
					\item[NIDNTTE] \phantom{.}
						\begin{compactdesc}
							\item[4:496] In \hyperref[HEBREWS-5-4]{Heb. 5:4} \gr{τιμή} denotes either the honor attached to a position or the position itself...We may say that \gr{τιμή} is regularly used to denote high valuation based on position in an organized whole. This whole is God's world of human beings, animals, and things. The distinctive distribution of honor among entities of varying worth is important. The resultant order and grades of honor must be respected not only by the one placed in a ``lower'' position, but also by the one placed ``higher.''
							\item[4:497] The honor of a human being and the honor of office are not mutually exclusive any more than are the honor of a human being and the honor of God.
							\item[4:497] The NT does not reduce the ground for honor, as did the Stoics, to an inner quality, such as virtue or wisdom.
								\begin{compactitem}
									\item See TDNTA 1072 above. I think in God we see the unification of both the inward and outward senses of ``honor'' here, which helps us make sense not only of the Hebrews use but also of DC 29.
								\end{compactitem}
						\end{compactdesc}
						
					% \item[LSJ , PDF ] \phantom{.}
						% \begin{compactitem}
							% \item 
						% \end{compactitem}
						
				\end{compactdesc}
				
				\begin{tabular}{ll}
					Total occurrences in the NT & 41/42
						% From my exhaustive concordance, the "41/42" is there, not sure what the disputed case is
				\end{tabular}
				
				\begin{compactdesc}
					\item[Some NT uses] \textbf{}
						\begin{compactdesc}
							\item[{\hyperref[HEBREWS-5-4]{Hebrews 5:4}}] \gr{καὶ οὐχ ἑαυτῷ τις λαμβάνει τὴν τιμήν, ἀλλὰ καλούμενος ὑπὸ τοῦ θεοῦ, καθώσπερ καὶ Ἀαρών.}
						\end{compactdesc}
				\end{compactdesc}
				
			% %
			% \subparagraph{TDNT: \texorpdfstring{\gr{sample word}}{} (3:536--556)}
			% %
				% \label{KATALLAGE-TDNT}
		
	% -----
	\subsubsection{Book of Mormon} % *BOM
	% -----
		\label{HONOR-BOM}
		
		% --
		\paragraph{Honor (noun)}
		% -- 
	
			\begin{tabular}{ll}
				Total occurrences in the BOM & 2
			\end{tabular}
			
			\begin{compactdesc}
				\item[{\hyperref[ALMA-1-16]{Alma 1:16}}]
				\item[{\hyperref[ALMA-60-36]{Alma 60:36}}]
			\end{compactdesc}
			
		% --
		\paragraph{Honor (verb)}
		% -- 
	
			\begin{tabular}{ll}
				Total occurrences in the BOM & 3
			\end{tabular}
			
			\begin{compactdesc}
				\item[{\hyperref[1NEPHI-17-55]{1 Nephi 17:55}}]
				\item[{\hyperref[2NEPHI-27-25]{2 Nephi 27:25}}]
				\item[{\hyperref[MOSIAH-13-20]{Mosiah 13:20}}]
			\end{compactdesc}
			
		% --
		\paragraph{Honorable}
		% -- 
	
			\begin{tabular}{ll}
				Total occurrences in the BOM & 3
			\end{tabular}
			
			Note that all three of these uses are quotations of Isaiah.
			
			\begin{compactdesc}
				\item[{\hyperref[2NEPHI-13-3]{2 Nephi 13:3}}]
				\item[{\hyperref[2NEPHI-13-5]{2 Nephi 13:5}}]
				\item[{\hyperref[2NEPHI-15-13]{2 Nephi 15:13}}]
			\end{compactdesc}
		
	% -----
	\subsubsection{Doctrine and Covenants} % *DC
	% -----
		\label{HONOR-DC}
	
		% \begin{tabular}{ll}
			% Total occurrences in the DC & 
		% \end{tabular}
		
		\begin{compactdesc}
			\item[{\hyperref[DC29-36]{DC 29:36}}]
			\item[{\hyperref[DC124-18]{DC 124:18}}] \phantom{.}
				\begin{compactitem}
					\item An unusual verse which talks about ``begetting'' glory and honor.
				\end{compactitem}
		\end{compactdesc}
		
	% -----
	\subsubsection{Pearl of Great Price} % *PGP
	% -----
		\label{HONOR-PGP}
	
		\begin{tabular}{ll}
			Total occurrences in the PGP & 1
		\end{tabular}
		
		\begin{compactdesc}
			\item[{\hyperref[MOSES-4-1]{Moses 4:1}}]
		\end{compactdesc}
		
% ----------------------
\subsection{Key Phrases} % *KEYPHRASES *PHRASES
% ----------------------
	\label{HONOR-PHRASES}

	\begin{compactdesc}
		\item[honor, power, glory] \textbf{7} | \hyperref[ALMA-60-36]{Alma 60:36}; \hyperref[DC20-36]{DC 20:36}; \hyperref[DC65-6]{DC 65:6}; \hyperref[DC84-102]{DC 84:102}; \hyperref[DC109-77]{DC 109:77}; \hyperref[DC128-21]{DC 128:21}; \hyperref[DC128-23]{DC 128:23}
			\begin{compactdesc}
				\item[Bible anchor verses] ---
					% \begin{compactdesc}
						% \item[Romans 5:5] \gr{}
					% \end{compactdesc}
				% \item[Commentary] ---
			\end{compactdesc}
	\end{compactdesc}
	
% % ----------------------
% \subsection{Bible Dictionaries} % *DICTIONARIES
% % ----------------------
	% \label{HONOR-BD}

% % ----------------------
% \subsection{Restoration Usage} % *RESTORATION
% % ----------------------
	% \label{HONOR-RESTORATION}

	% % -----
	% \subsubsection{Joseph Smith} % *JS
	% % -----
		% \label{HONOR-JS}
	
		% \begin{compactdesc}
			% \item[{\hyperref[KEY]{Date/Source}}]
		% \end{compactdesc}
		
	% % -----
	% \subsubsection{Temple Ordinances}
	% % -----
		% \label{HONOR-TEMPLE}
	
		% \begin{compactdesc}
			% \item[{\hyperref[KEY]{Date/Source}}]
		% \end{compactdesc}
	
	% % -----
	% \subsubsection{Brigham Young} % *BY
	% % -----
		% \label{HONOR-BY}
	
		% \begin{compactdesc}
			% \item[{\hyperref[KEY]{Date/Source}}]
		% \end{compactdesc}
	
% % ----------------------
% \subsection{Commentary and Studies} % *COMMENTARY *STUDIES
% % ----------------------
	% \label{HONOR-COMMENTARY}

	% % -----
	% \subsubsection{Key Points}
	% % -----
		% \label{HONOR-KEYS}
	
		% \begin{compactitem}
			% \item
		% \end{compactitem}
		
	% % -----
	% \subsubsection{Sample Study}
	% % -----
		% \label{HONOR-study}